\appendix

\section{The Model in Detail}\label{app:model}

The generic model puts the pool inside a stateful wrapper, \m{reactive_system}:
\begin{lstlisting}
model runtime init_policy service_policy
      restore_invariant p_pub p_priv p_sched =
  reactive_system init_policy
    (state_step service_policy restore_invariant)
    no_output
    (pool runtime p_pub p_priv p_sched)
    pool_input
  : Proc T_in T_out
\end{lstlisting}
\m{no_output}, the all-\code{None} tuple, stands in for the pool's output before its
first step, and \m{pool_input} passes each slot the current pid and the disk
handler's output.

\paragraph{The state} Besides the current pid and the interrupt controller, the
state holds the pid of the user process a handler interrupted, a reschedule flag
\m{re_sched}, and a time-slice counter \m{ir_count}. The last two are used only by
the sliced policy.

\paragraph{Handlers} \code{ir_handler runtime} counts its output steps modulo
\code{runtime}. Its output type \code{THandlerOutput} has two values: it emits
\code{Notify} on the step that completes a run and \code{Nothing} (a handler value,
not the option \code{None}) otherwise.

\paragraph{The pool} Each slot is a \code{swi} whose gate bit is computed from the
pid. Each slot's output is tagged with the constant \code{true}, which closes its
switch after one output step, so a selected process takes exactly one step.

\paragraph{The scheduling decision} The state update is
\begin{lstlisting}
state_step service_policy restore_invariant =
  step_sum set_pending
    (initiate_next restore_invariant
     (@$\circ$@) service_policy (@$\circ$@) check_scheduler)
\end{lstlisting}
where \code{step_sum} applies its first argument to interrupts and its second to
pool outputs, and the composition runs right to left. The priority rule of
\cref{sec:state} is
\begin{lstlisting}
initiate_next restore_invariant v =
  if masks_set v then v   (* handler running *)
  else let v := restore_invariant v in
    if first_ready v is Some ir
    then initiate_handler ir v
    else if is_handler_pid v (* handler done *)
    then if get_re_sched v
         then initiate_scheduler v
         else initiate_prev_pid v
    else v                (* stay in user space *)
\end{lstlisting}
\m{first_ready} picks the highest-priority serviceable interrupt.
\m{initiate_handler} saves the current user pid, points the current pid at the
handler, sets all masks and clears the handler's pending bit.
\m{initiate_scheduler} hands control to the scheduler and clears \m{re_sched};
\m{initiate_prev_pid} returns to the interrupted user process.
